%--------------------------------------------------------------------------------------------------
% 
\chapter*{Abstract}
\pdfbookmark[0]{Abstract}{Abstract}
%--------------------------------------------------------------------------------------------------

The ever-increasing amount of global news presents an overwhelming challenge for individuals and organizations striving to keep abreast of relevant information. 
The sheer volume of unstructured text data makes it nearly impossible for users to consume and comprehend all the information pertinent to their interests.
Media monitoring and analysis systems can help us by employing advanced algorithms to address the problem of size and the problems of opinionated views and misinformation.
These tools allow users to quickly glean the essential information from an article, saving time and reducing cognitive load.
This paper presents an overview of two specific topics concerning news analysis and Natural Language Processing (NLP): Named Entity Recognition and News Framing.
In the context of these NLP tasks, we recognize the potential for utilizing cross-lingual NLP techniques to augment the analysis of news in languages that lack sufficient annotated data.
The primary goal of cross-lingual NLP is to transfer knowledge from high-resource languages to low-resource languages where annotated data is scarce or unavailable.
We mainly focus on low-resource Slovene and other South Slavic languages with even less annotated data.
Therefore we analyze key techniques in the needed cross-lingual setting, which include multilingual pre-training, transfer learning, and zero-shot learning.
